% Source: https://artofproblemsolving.com/community/c3h1772117_please_help
% Stable identity: aops-topic-1772117-q01
Let the barrel hold $V$ litres, so it contained $\dfrac{2V}{3}$ litres. Working backwards is easiest.

\textbf{Second draw.} Taking out $1$ litre less than a third of the amount $R$ leaves $\dfrac{2}{3}R+1$ litres. So $\dfrac{2}{3}R+1=15$, $\dfrac{2}{3}R=14$ and $R=21$.

\textbf{First draw.} Taking out $1$ litre less than half of the contents $\dfrac{2V}{3}$ leaves half of it plus $1$: $\dfrac{V}{3}+1=21$, so $V=60$.

The barrel held $\boxed{60}$ litres. Check: $40$ litres of wine; he took $20-1=19$, leaving $21$; then $7-1=6$, leaving $15$.
